\documentclass[11pt]{article}
\usepackage[a4paper,margin=1in]{geometry}
\usepackage{amsmath,amssymb,amsthm,booktabs,enumitem}
\usepackage[T1]{fontenc}
\usepackage{lmodern}
\usepackage{hyperref}

\newtheorem{theorem}{Theorem}[section]
\newtheorem{proposition}[theorem]{Proposition}
\newtheorem{lemma}[theorem]{Lemma}
\newtheorem{corollary}[theorem]{Corollary}
\theoremstyle{definition}
\newtheorem{definition}[theorem]{Definition}
\theoremstyle{remark}
\newtheorem{remark}[theorem]{Remark}
\newcommand{\Mob}{\mu}
\newcommand{\ellobs}{\ell}
\newcommand{\abs}[1]{\lvert#1\rvert}

\title{Bounded constraint-graph transducer certificates\\
for arithmetic capacity}
\author{RH research program}
\date{August 2026}

\begin{document}
\maketitle

\begin{abstract}
We isolate a finite, reusable theorem behind several arithmetic capacity
experiments.  A finite directed constraint graph with a bounded vertex
observable admits an exact open-prefix max-plus dynamic program.  A clocked
finite-memory transducer is called universally safe when every possible pair
of inputs produces a legal graph edge.  If its observed label is a one-site
factor of the current clock phase and the current Mobius value, then its
Mobius correlation has an unconditional arithmetic-progression formula in
terms of squarefree densities.  Consequently every such certificate gives a
rigorous lower bound on the graph capacity liminf, and all certificates with
fixed clock and memory budgets form a finite enumerable class.  We verify the
contract on the RH-366 four-state graph, the distinct RH-368 three-cell
factor, and a new period-three safe switch on the RH-366 graph.  Labels that
depend on memory are explicitly left open because they require higher-order
Mobius correlations.  This is a bounded classification result, not a
classification of all subshifts and not a spectral or zeta construction.
\end{abstract}

\section{Frozen data and claim boundary}

RH-366 fixes a four-state survivor whose state order is
\((--,-+,+-,++)\) and whose adjacency matrix is
\begin{equation}
A_{366}=\begin{pmatrix}
1&0&1&0\\
1&0&0&0\\
0&1&0&1\\
0&1&0&0
\end{pmatrix},
\qquad
\ell_{366}=(-1,-1,1,1).
\label{eq:rh366graph}
\end{equation}
The graph path records a pair of consecutive signs, so this is the
distance-two constraint used by RH-366 \cite{RH366}.  RH-368 supplies a
different three-cell graph
\begin{equation}
A_{368}=\begin{pmatrix}0&0&1\\0&0&1\\1&1&0\end{pmatrix},
\qquad
\ell_{368}=(1,-1,-1),
\label{eq:rh368graph}
\end{equation}
with its parity factor and all-order capacity theorem \cite{RH368}.  We use
both only as frozen finite graph inputs; the two languages are not identified.

For a finite directed graph \(G=(V,E)\) and an integer observable
\(\ellobs:V\to\mathbb Z\), write
\begin{equation}
 K_N(G,\ellobs)=
 \max_{\substack{v_1,\ldots,v_N\in V\\(v_n,v_{n+1})\in E}}
 \abs{\sum_{n=1}^N \Mob(n)\ellobs(v_n)}.
 \label{eq:capacity}
\end{equation}
The optimizing open path may depend on the complete prefix.  This definition
is a finite scalar functional; it is not a trace or a determinant.

\section{Exact graph capacity}

\begin{theorem}[Open-path max-plus recurrence]
Assume every vertex has an incoming edge (as in the frozen graphs above).  Put
\begin{equation}
D_1^{\pm}(v)=\pm\Mob(1)\ellobs(v),
\qquad
D_{n+1}^{\pm}(v)=\pm\Mob(n+1)\ellobs(v)
 +\max_{(u,v)\in E}D_n^{\pm}(u).
\label{eq:dp}
\end{equation}
Then
\begin{equation}
K_N(G,\ellobs)=\max_{v\in V}\max\{D_N^+(v),D_N^-(v)\}.
\label{eq:dp-capacity}
\end{equation}
The recurrence and a maximizing witness require \(O(N\abs{E})\) integer
operations.
\end{theorem}

\begin{proof}
Induct on \(n\).  The value \(D_n^+(v)\) is the largest signed score among
paths ending at \(v\): appending \(v\) contributes
\(\Mob(n)\ellobs(v)\), and the predecessor must be an incoming neighbor.
The same induction with the opposite sign gives \(D_n^-(v)\), the largest
negative of a path score.  Taking the maximum over terminal vertices proves
\eqref{eq:dp-capacity}.  Each edge is inspected once per step; predecessor
pointers give a witness.
\end{proof}

\begin{corollary}[Universal squarefree upper bound]
For bounded \(\ellobs\),
\begin{equation}
\limsup_{N\to\infty}\frac{K_N(G,\ellobs)}N
\leq \frac{6}{\pi^2}\|\ellobs\|_\infty.
\label{eq:upper}
\end{equation}
\end{corollary}

\begin{proof}
Every admissible path has absolute score at most
\(\|\ellobs\|_\infty\sum_{n\leq N}\Mob(n)^2\).  The squarefree density is
\(6/\pi^2\).
\end{proof}

\section{Safe finite-memory transducers}

\begin{definition}[Clocked transducer]
Fix a clock \(q\geq1\), a finite state set \(S\), and input alphabet
\(\mathcal A=\{-1,0,1\}\).  A transducer consists of maps
\[
 \tau:S\times\mathbb Z/q\mathbb Z\times\mathcal A\to S,
 \qquad
 \omega:S\times\mathbb Z/q\mathbb Z\times\mathcal A\to V.
\]
At time \(n\), with \(a_n=\Mob(n)\) and \(r_n=n\bmod q\), it outputs
\(v_n=\omega(s_n,r_n,a_n)\) and updates
\(s_{n+1}=\tau(s_n,r_n,a_n)\).
\end{definition}

\begin{definition}[Universal safety and one-site factor]
The table \((\tau,\omega)\) is universally safe on \(G\) if, for every
\(s\in S\), \(r\in\mathbb Z/q\mathbb Z\), and \(a,b\in\mathcal A\),
\begin{equation}
\bigl(\omega(s,r,a),
 \omega(\tau(s,r,a),r+1,b)\bigr)\in E.
\label{eq:safety}
\end{equation}
It has a one-site observed factor if there are functions
\(g_r:\mathcal A\to\mathbb Z\) with
\begin{equation}
 \ellobs(\omega(s,r,a))=g_r(a)
 \quad\text{for every }s,r,a.
\label{eq:onesite}
\end{equation}
\end{definition}

The universal quantifier in \eqref{eq:safety} is intentional.  It is a
finite table check and does not assume anything about future Mobius values.
The one-site condition is the arithmetic firewall: the internal memory may
choose a legal path, but it may not alter the observed label.

\begin{lemma}[Squarefree density in a residue class]
For \(r\in\mathbb Z/q\mathbb Z\),
\begin{equation}
 \frac1N\#\{n\leq N:n\equiv r\pmod q,\Mob(n)^2=1\}
 \longrightarrow
 \delta_{q,r}:=
 \sum_{\substack{d\geq1\\(q,d^2)\mid r}}
 \frac{\Mob(d)}{\operatorname{lcm}(q,d^2)}.
\label{eq:delta}
\end{equation}
\end{lemma}

\begin{proof}
Use \(\Mob(n)^2=\sum_{d^2\mid n}\Mob(d)\).  The simultaneous congruences
\(n\equiv r\pmod q\) and \(d^2\mid n\) are soluble exactly when
\((q,d^2)\mid r\), and then have density
\(1/\operatorname{lcm}(q,d^2)\).  Truncating at \(d\leq\sqrt N\) gives an
\(O(\sqrt N)\) counting error; the tail is \(O(N/\sqrt N)\).  Absolute
convergence of the displayed series completes the limit.
\end{proof}

\begin{theorem}[One-site transducer correlation]
Let a universally safe transducer satisfy \eqref{eq:onesite}.  Then its
explicit Mobius path obeys
\begin{equation}
 \lim_{N\to\infty}\frac1N\sum_{n\leq N}\Mob(n)\ellobs(v_n)
 =L(T):=\frac12\sum_{r\bmod q}\delta_{q,r}
       \bigl(g_r(1)-g_r(-1)\bigr).
\label{eq:limit}
\end{equation}
Consequently,
\begin{equation}
 |L(T)|\leq\liminf_{N\to\infty}\frac{K_N(G,\ellobs)}N
 \leq\limsup_{N\to\infty}\frac{K_N(G,\ellobs)}N
 \leq\frac6{\pi^2}\|\ellobs\|_\infty.
\label{eq:bracket}
\end{equation}
\end{theorem}

\begin{proof}
For each residue \(r\), let \(Q_r(N)\) be the squarefree count in
\eqref{eq:delta} and put
\[
 M_r(N)=\sum_{\substack{n\leq N\\n\equiv r\pmod q}}\Mob(n).
\]
Davenport's
fixed-frequency estimate, followed by finite Fourier inversion, gives
\(M_r(N)=o(N)\).  Hence the positive and negative counts in that residue are
\((Q_r\pm M_r)/2\).  Since zero Mobius values contribute zero to the
correlation, the residue-\(r\) contribution divided by \(N\) is
\[
 \frac{Q_r}{2N}(g_r(1)-g_r(-1))
 +\frac{M_r}{2N}(g_r(1)+g_r(-1)).
\]
Apply the lemma and sum over the finitely many residues.  The explicit path
is admissible by \eqref{eq:safety}, so its absolute score is bounded by the
capacity; the upper bound is \eqref{eq:upper}.
\end{proof}

\begin{proposition}[Finite bounded-resource classification]
Fix \(G,q\), and a memory budget \(m=\abs S\).  The set of all transducer
tables satisfying universal safety and the one-site condition is finite and
can be exhaustively enumerated.  If
\[
 \Gamma_{q,m}(G,\ellobs)=\max_T\abs{L(T)},
\]
where the maximum is over that finite set, then
\[
 \Gamma_{q,m}(G,\ellobs)\leq\liminf_{N\to\infty}K_N(G,\ellobs)/N.
\]
\end{proposition}

\begin{proof}
There are finitely many choices for the transition and output entries of
the tables \(S\times(\mathbb Z/q\mathbb Z)\times\mathcal A\), and finitely
many initial states.  Conditions \eqref{eq:safety} and \eqref{eq:onesite}
are finite Boolean tests.  The inequality follows by maximizing the
individual lower bounds in \eqref{eq:bracket}.
\end{proof}

\begin{remark}[What the finite classification does not say]
The proposition is not a classification of all mixing subshifts: \(q\) and
the memory budget are fixed before enumeration.  If the observed label
depends on \(s_n\), the residue reduction is invalid and higher-order Mobius
correlations enter.  We do not replace those missing correlations by pair
data or by a finite numerical fit.
\end{remark}

\section{Three frozen certificates}

\subsection{The RH-366 four-state label rule}
Use the two-state universal safety completion supplied in the artifact.  It
has the same one-site labels as the RH-366 rule.  Set
\[
 e_n=+1\quad\Longleftrightarrow\quad
 \Mob(n)=1\text{ and }n\bmod4\in\{1,2\},
\]
and output a graph representative with that label.  The state-1 rows at
phases 2 and 3 are completion rows; they are not an additional arithmetic
observable.  The selected residues shift by two to \(3,0\pmod4\), so two
selected plus signs can never be two places apart.  The completed state table
is universally safe and its observed factor has
\(g_1(1)-g_1(-1)=g_2(1)-g_2(-1)=2\), with the other differences zero.
Thus
\[
 L(T)=\delta_{4,1}+\delta_{4,2}=\frac4{\pi^2}.
\]
This recovers the RH-366 lower certificate, not its unresolved capacity
limit.

\subsection{The RH-368 three-cell factor}
For \eqref{eq:rh368graph}, at odd phases output vertex \(0\) when
\(\Mob(n)=1\) and vertex \(1\) otherwise; at even phases output vertex
2.  Every output edge is legal.  Since \(\ellobs=(1,-1,-1)\), this has
\(g_1(1)-g_1(-1)=2\) and gives
\[
 L(T)=\delta_{2,1}=\frac4{\pi^2}.
\]
RH-368 proves the separate parity-factor capacity limit; here it is used as
an independent frozen graph audit of the general certificate contract.

\subsection{A new period-three safe switch}
On the RH-366 graph, use the anchor state \(0=(--)\).  At a phase
\(n\equiv1\pmod3\), choose the loop
\[
 0\to2\to1\to0
\]
when \(\Mob(n)=1\), and choose
\[
 0\to0\to0\to0
\]
otherwise.  The two paths have equal length and their observable labels
\((-1,1,-1,-1)\) versus \((-1,-1,-1,-1)\) differ only at the first output
position.  The finite state table also defines safe outputs from every
unreachable state, so the universal check is literal rather than
trajectory-dependent.  Therefore
\[
 g_1(1)-g_1(-1)=2,
 \qquad L(T)=\delta_{3,1}=\frac9{4\pi^2}.
\]
This is a new bounded certificate, not a claim about the RH-366 optimizer's
limit.

\begin{table}[t]
\centering
\small
\begin{tabular}{@{}lllll@{}}
\toprule
Instance & graph vertices & clock & memory & certified limit \\
\midrule
RH-366 rule & 4 & 4 & 2 & $4/\pi^2$ \\
RH-368 factor & 3 & 2 & 1 & $4/\pi^2$ \\
RH-366 q=3 switch & 4 & 3 & 2 & $9/(4\pi^2)$ \\
\bottomrule
\end{tabular}
\caption{Finite safe transducer certificates.  The constants are exact
one-site limits; they are lower bounds for the corresponding open capacities.}
\label{tab:instances}
\end{table}

\section{Executable audit}

The artifact uses a linear integer Mobius sieve, exact max-plus recurrences,
literal universal table checks, and a small exhaustive table enumeration.
It checks \(N\leq128\) for all three transducer witnesses, an endpoint
\(N=2^{16}\), and all \(3^6=729\) one-state output tables on the RH-368
graph at \(q=2\).  The source manifest locks nine files and five source
commits.  These are finite reproduction and provenance checks.  The only
asymptotic steps are the squarefree-density lemma and the cited Davenport
estimate used in the proof.

\section{Route verdict and Gate ledger}

Route A is \texttt{GO} narrowly: the max-plus theorem, the finite safety
classification, the one-site arithmetic limit, and the three audited
certificates form an independent theorem package.  Route B is
\texttt{STOP\_SCOPED}.  The transducer is an offline arithmetic selector, so
the first mismatch occurs before Gate A: no canonical intrinsic operator,
determinant, signed von-Mangoldt trace, or completed-zeta divisor is present.
The RH-366 capacity limit remains open; memory-dependent labels remain
outside the unconditional theorem.

The five project Gates remain false/open.  This paper does not construct a
Hilbert--Polya operator, identify Riemann zeros, prove a prime-power trace, or
prove the Riemann Hypothesis.

\section{Conclusion}

Finite constraint graphs admit a precise arithmetic certificate layer: exact
finite optimization, a decidable bounded transducer contract, and an
unconditional one-site Mobius limit.  The bounded qualifier is the result's
strength and its boundary.  Increasing memory or allowing the observable to
read that memory is a new higher-order arithmetic problem, not something a
finite table or pair ledger can silently solve.

\bibliographystyle{plain}
\bibliography{references}

\end{document}
